% GENERATED FILE. Edit canonical lessons, then run npm run generate:books.
% canonical-source-hash: fnv1a64:da4b93ec8478135a
% canonical-lessons: GU-C417-pol, GU-C417-deli, GU-C417-otlo, GU-C417-khadki, GU-C417-vasvu-shut

\chapter{The Pol}
\label{ch:gu-the-pol}
\hlchaptermodality{417}

\begin{chapteropening}
\textbf{By the end of this chapter:} \emph{I can describe an old-city pol and its gateway.}

Complete the last lesson of chapter 417: I can describe an old-city pol and its gateway.
\end{chapteropening}

\section[\texorpdfstring{\gu{પોળ} (poḷ)}{પોળ (poḷ)}]{\gu{પોળ} (poḷ) — a pol (a gated old-city neighbourhood)}

\label{lesson:GU-C417-pol}



\begin{quote}
Before the new one: say the Gujarati for the outskirts of a village, then the Gujarati for to travel about.
\end{quote}

\subsection*{You'll want to know: \gu{પોળ}}
\textbf{\gu{પોળ}} — \emph{poḷ} — \textquotedblleft{}a pol (a gated old-city neighbourhood)\textquotedblright{}.

\begin{quote}
\gu{જૂના} \gu{શહેરની} \gu{પોળમાં} \gu{સાંકડી} \gu{ગલીઓ} \gu{છે}.
\end{quote}

\begin{quote}
\emph{jūnā śahernī poḷmā̃ sā̃kḍī galīo che.}
\end{quote}

\begin{quote}
The pol in the old city has narrow lanes.
\end{quote}

A feminine noun.

\begin{cousinweb}
\textbf{\gu{પોળ}} (\emph{poḷ}) is a tadbhava of Sanskrit \emph{pratōlī}, \textquotedblleft{}a main street; a town gate\textquotedblright{}; CDIAL lists the Gujarati word, \textquotedblleft{}a street\textquotedblright{} (CDIAL 8633). A \gu{પોળ} is a street of houses shut in by its own gate.
\end{cousinweb}

\begin{grammarlens}[title={What you've built}]
One more word for getting around.
\end{grammarlens}

\subsection*{Your turn}
\begin{itemize}
\raggedright
  \item \emph{Say it:} \emph{poḷ}
  \item \emph{Say it:} \emph{poḷ}, once more
  \item \emph{Say it:} \emph{jūnā śahernī poḷmā̃ sā̃kḍī galīo che.} — The pol in the old city has narrow lanes.
  \item \emph{Recall:} say the Gujarati for the outskirts of a village, then the Gujarati for to travel about, then say \emph{poḷ} again
\end{itemize}

\begin{tcolorbox}[breakable,colback=teal!4,colframe=teal!35!black,title={Before you move on}]
What does \gu{પોળ} mean? (\textquotedblleft{}A pol (a gated old-city neighbourhood)\textquotedblright{}.) Say it once more.
\end{tcolorbox}
\section[\texorpdfstring{\gu{ડેલી} (ḍelī)}{ડેલી (ḍelī)}]{\gu{ડેલી} (ḍelī) — an entrance gateway (of a house or a pol)}

\label{lesson:GU-C417-deli}



\begin{quote}
Before the new one: say the Gujarati for to travel about, then the Gujarati for a pol.
\end{quote}

\subsection*{You'll want to know: \gu{ડેલી}}
\textbf{\gu{ડેલી}} — \emph{ḍelī} — \textquotedblleft{}an entrance gateway (of a house or a pol)\textquotedblright{}.

\begin{quote}
\gu{ડેલી} \gu{પાસે} \gu{ચોકીદાર} \gu{બેઠો} \gu{છે}.
\end{quote}

\begin{quote}
\emph{ḍelī pāse cokīdār beṭho che.}
\end{quote}

\begin{quote}
A watchman is sitting by the gateway.
\end{quote}

A feminine noun.

\begin{cousinweb}
\textbf{\gu{ડેલી}} (\emph{ḍelī}) is a tadbhava of Sanskrit \emph{dēhalī}, \textquotedblleft{}a threshold\textquotedblright{}; CDIAL lists the Gujarati word for a porch or a covered entrance (CDIAL 6559).
\end{cousinweb}

\begin{grammarlens}[title={What you've built}]
One more word for getting around.
\end{grammarlens}

\subsection*{Your turn}
\begin{itemize}
\raggedright
  \item \emph{Say it:} \emph{ḍelī}
  \item \emph{Say it:} \emph{ḍelī}, once more
  \item \emph{Say it:} \emph{ḍelī pāse cokīdār beṭho che.} — A watchman is sitting by the gateway.
  \item \emph{Recall:} say the Gujarati for to travel about, then the Gujarati for a pol, then say \emph{ḍelī} again
\end{itemize}

\begin{tcolorbox}[breakable,colback=teal!4,colframe=teal!35!black,title={Before you move on}]
What does \gu{ડેલી} mean? (\textquotedblleft{}An entrance gateway (of a house or a pol)\textquotedblright{}.) Say it once more.
\end{tcolorbox}
\section[\texorpdfstring{\gu{ઓટલો} (oṭlo)}{ઓટલો (oṭlo)}]{\gu{ઓટલો} (oṭlo) — a raised platform (at the front of a house)}

\label{lesson:GU-C417-otlo}



\begin{quote}
Before the new one: say the Gujarati for a pol, then the Gujarati for an entrance gateway.
\end{quote}

\subsection*{You'll want to know: \gu{ઓટલો}}
\textbf{\gu{ઓટલો}} — \emph{oṭlo} — \textquotedblleft{}a raised platform (at the front of a house)\textquotedblright{}.

\begin{quote}
\gu{સાંજે} \gu{દાદી} \gu{ઓટલા} \gu{પર} \gu{બેસે} \gu{છે}.
\end{quote}

\begin{quote}
\emph{sā̃je dādī oṭlā par bese che.}
\end{quote}

\begin{quote}
In the evening grandmother sits on the front platform.
\end{quote}

A masculine noun.

\begin{cousinweb}
\textbf{\gu{ઓટલો}} (\emph{oṭlo}) goes back to \emph{*ōṭṭā}, a form CDIAL reconstructs; CDIAL lists the Gujarati word, a raised platform to sit on (CDIAL 2544).
\end{cousinweb}

\begin{grammarlens}[title={What you've built}]
One more word for getting around.
\end{grammarlens}

\subsection*{Your turn}
\begin{itemize}
\raggedright
  \item \emph{Say it:} \emph{oṭlo}
  \item \emph{Say it:} \emph{oṭlo}, once more
  \item \emph{Say it:} \emph{sā̃je dādī oṭlā par bese che.} — In the evening grandmother sits on the front platform.
  \item \emph{Recall:} say the Gujarati for a pol, then the Gujarati for an entrance gateway, then say \emph{oṭlo} again
\end{itemize}

\begin{tcolorbox}[breakable,colback=teal!4,colframe=teal!35!black,title={Before you move on}]
What does \gu{ઓટલો} mean? (\textquotedblleft{}A raised platform (at the front of a house)\textquotedblright{}.) Say it once more.
\end{tcolorbox}
\section[\texorpdfstring{\gu{ખડકી} (khaḍkī)}{ખડકી (khaḍkī)}]{\gu{ખડકી} (khaḍkī) — a small gated lane (inside a pol)}

\label{lesson:GU-C417-khadki}



\begin{quote}
Before the new one: say the Gujarati for an entrance gateway, then the Gujarati for a raised platform.
\end{quote}

\subsection*{You'll want to know: \gu{ખડકી}}
\textbf{\gu{ખડકી}} — \emph{khaḍkī} — \textquotedblleft{}a small gated lane (inside a pol)\textquotedblright{}.

\begin{quote}
\gu{અમારી} \gu{ખડકીમાં} \gu{પાંચ} \gu{ઘર} \gu{છે}.
\end{quote}

\begin{quote}
\emph{amārī khaḍkīmā̃ pā̃c ghar che.}
\end{quote}

\begin{quote}
There are five houses in our little lane.
\end{quote}

A feminine noun.

\begin{cousinweb}
\textbf{\gu{ખડકી}} (\emph{khaḍkī}) is a tadbhava of Sanskrit \emph{khaṭakkikā}, \textquotedblleft{}a side door\textquotedblright{}; CDIAL lists the Gujarati word, \textquotedblleft{}a courtyard, a lane\textquotedblright{} (CDIAL 3770). The name of the small gate seems to have passed to the lane behind it.
\end{cousinweb}

\begin{grammarlens}[title={What you've built}]
One more word for getting around.
\end{grammarlens}

\subsection*{Your turn}
\begin{itemize}
\raggedright
  \item \emph{Say it:} \emph{khaḍkī}
  \item \emph{Say it:} \emph{khaḍkī}, once more
  \item \emph{Say it:} \emph{amārī khaḍkīmā̃ pā̃c ghar che.} — There are five houses in our little lane.
  \item \emph{Recall:} say the Gujarati for an entrance gateway, then the Gujarati for a raised platform, then say \emph{khaḍkī} again
\end{itemize}

\begin{tcolorbox}[breakable,colback=teal!4,colframe=teal!35!black,title={Before you move on}]
What does \gu{ખડકી} mean? (\textquotedblleft{}A small gated lane (inside a pol)\textquotedblright{}.) Say it once more.
\end{tcolorbox}
\section[\texorpdfstring{\gu{વાસવું} (vāsvũ)}{વાસવું (vāsvũ)}]{\gu{વાસવું} (vāsvũ) — to shut (a door)}

\label{lesson:GU-C417-vasvu-shut}



\begin{quote}
Before the new one: say the Gujarati for a raised platform, then the Gujarati for a small gated lane.
\end{quote}

\subsection*{You'll want to know: \gu{વાસવું}}
\textbf{\gu{વાસવું}} — \emph{vāsvũ} — \textquotedblleft{}to shut (a door)\textquotedblright{}.

\begin{quote}
\gu{રાત્રે} \gu{દાદાએ} \gu{ડેલીનું} \gu{બારણું} \gu{વાસ્યું}.
\end{quote}

\begin{quote}
\emph{rātre dādāe ḍelīnũ bārṇũ vāsyũ.}
\end{quote}

\begin{quote}
At night grandfather shut the door of the gateway.
\end{quote}

Past, with \gu{તેણે} (he or she …): \gu{તેણે} \gu{બારણું} \gu{વાસ્યું}, \emph{teṇe bārṇũ vāsyũ}.

\begin{cousinweb}
\textbf{\gu{વાસવું}} (\emph{vāsvũ}) No dictionary we use records an older form, so its history is not settled.
\end{cousinweb}

\begin{grammarlens}[title={What you've built}]
One more word for getting around.
\end{grammarlens}

\subsection*{Your turn}
\begin{itemize}
\raggedright
  \item \emph{Say it:} \emph{vāsvũ}
  \item \emph{Say it:} \emph{vāsvũ}, once more
  \item \emph{Say it:} \emph{rātre dādāe ḍelīnũ bārṇũ vāsyũ.} — At night grandfather shut the door of the gateway.
  \item \emph{Recall:} say the Gujarati for a raised platform, then the Gujarati for a small gated lane, then say \emph{vāsvũ} again
\end{itemize}

\begin{tcolorbox}[breakable,colback=teal!4,colframe=teal!35!black,title={Before you move on}]
What does \gu{વાસવું} mean? (\textquotedblleft{}To shut (a door)\textquotedblright{}.) Say it once more.
\end{tcolorbox}
